\chapter{Statistics for Data Scientists: Sampling, CLT, Confidence Intervals and Hypothesis Testing}
\companion{StatQuest: p-values, what they are and how to interpret them}{\ytid{vemZtEM63GY}}

The checklist recommends the Rice University statistics courses and lists descriptive vs inferential statistics, population vs sample, the central limit
theorem, p-values and hypothesis testing. The reported Round 1 question was exactly ``what is a p-value, what is alpha, how do you compare them''. This
chapter teaches inference from zero with worked numbers, and then shows where it drives ML decisions (model comparison, A/B tests, drift).

\section{Descriptive vs inferential statistics}
\begin{itemize}
\item \textbf{Descriptive statistics} summarise the data you have: mean, median, standard deviation, histograms (Chapter 31). ``The 200 applicants to this job had a
  median experience of 3 years.''
\item \textbf{Inferential statistics} use a sample to draw conclusions about a larger population, with quantified uncertainty: estimates, confidence intervals,
  hypothesis tests. ``Seekers who receive the new alert apply 12\% more (95\% CI: 5\% to 19\%).''
\end{itemize}

\section{Population, sample, parameter, statistic}
\begin{itemize}
\item \textbf{Population}: everyone or everything you care about (all Naukri seekers in India). Its numerical characteristics are \textbf{parameters} ($\mu$, $\sigma$, $p$):
  fixed but usually unknown.
\item \textbf{Sample}: the subset you observe (10{,}000 surveyed seekers). Its numbers are \textbf{statistics} ($\bar x$, $s$, $\hat p$): known but random (a different sample gives
  different values).
\end{itemize}
\textbf{Sampling methods}: simple random; \textbf{stratified} (sample within groups such as cities, then combine: guarantees coverage of small groups, as in stratified CV);
cluster (sample whole groups such as colleges); systematic (every $k$-th record). \textbf{Sampling bias}: the sample differs systematically from the population (survey only
active users; train only on users who were called by telecallers). No sample size fixes bias.

\section{Sampling distributions and the standard error}
If we repeated the sampling many times, the statistic (say $\bar x$) would vary. Its distribution is the \textbf{sampling distribution}; its standard deviation is the
\textbf{standard error} (SE). For a mean of $n$ independent observations:
\[ \text{SE}(\bar x) = \frac{\sigma}{\sqrt n}\quad(\text{estimated by } s/\sqrt n). \]
For a proportion: $\text{SE}(\hat p) = \sqrt{p(1-p)/n}$. Four times the data halves the error.
\begin{example}[: standard error]
Test accuracy 0.80 on 400 examples: SE $= \sqrt{0.8\cdot0.2/400} = 0.02$. On 43 examples (Chapter 27): SE $= 0.061$: the estimate could easily be off by 10 points.
\end{example}

\section{The central limit theorem}
\textbf{Statement}: for independent observations from \emph{any} distribution with finite mean $\mu$ and variance $\sigma^2$, the sample mean is approximately normal for large $n$:
\[ \bar X\approx\mathcal N\Big(\mu,\frac{\sigma^2}{n}\Big), \]
however skewed the original distribution is. ``Large'' is often $n\ge30$; more for very skewed data.
\begin{example}[: CLT with dice]
One die is uniform on 1--6 (mean 3.5, variance 2.917). The average of 30 rolls is approximately normal with mean 3.5 and SE $\sqrt{2.917/30} = 0.31$. So about 95\% of such averages
fall within $3.5\pm0.61$, i.e. 2.89 to 4.11, even though a single roll is anything but normal.
\end{example}
\textbf{Why it matters}: it justifies normal-based confidence intervals and z/t-tests for means and proportions (A/B tests on clicks), explains why averages are stable,
and why sums of many small effects look Gaussian.
Related: the \textbf{law of large numbers} (the sample mean converges to $\mu$); the CLT adds \emph{how fast} and \emph{in what shape}.

\section{Confidence intervals}
A 95\% confidence interval is a range computed from the sample by a procedure that captures the true parameter in 95\% of repeated samples:
\[ \bar x\pm z^*\frac{s}{\sqrt n}\ \ (z^* = 1.96\text{ for }95\%),\qquad\text{or } \bar x\pm t^*_{n-1}\frac{s}{\sqrt n}\text{ for small }n. \]
\begin{example}[: CI for a mean]
Time-to-hire for 36 jobs: $\bar x = 12$ days, $s = 3$. SE $= 0.5$. 95\% CI $= 12\pm1.96\times0.5 = (11.02, 12.98)$ (with $t^*_{35} = 2.03$: $12\pm1.02$).
\end{example}
\textbf{Correct interpretation}: the method has 95\% coverage; it is \emph{not} ``95\% probability that $\mu$ is in this particular interval'' (in frequentist terms $\mu$ is fixed;
a Bayesian credible interval does have that reading). Wider intervals for higher confidence or smaller samples.

\section{Hypothesis testing, step by step}
\begin{enumerate}
\item \textbf{Null hypothesis $H_0$}: the default, ``no effect'' claim (the new ranking does not change apply rate; the coin is fair).
\item \textbf{Alternative $H_1$}: what we want evidence for (apply rate changed: two-sided; increased: one-sided).
\item \textbf{Significance level $\alpha$}: the false-positive rate we accept, chosen before seeing data (0.05).
\item \textbf{Test statistic}: how far the data are from $H_0$ in standard-error units (z, t, $\chi^2$, F).
\item \textbf{p-value}: the probability, \emph{assuming $H_0$ is true}, of getting a result at least as extreme as the one observed.
\item \textbf{Decision}: if p $\le\alpha$, reject $H_0$ (``statistically significant''); otherwise fail to reject (not the same as proving $H_0$).
\end{enumerate}
\begin{example}[: a z-test]
Average applications per job are claimed to be 10 (with known $\sigma = 2$). A sample of 64 jobs averages 10.5. $z = \frac{10.5 - 10}{2/\sqrt{64}} = \frac{0.5}{0.25} = 2$.
Two-sided p-value $= P(|Z|\ge2) = 0.046 < 0.05$: reject $H_0$ at the 5\% level (but not at 1\%).
\end{example}

\subsection{What a p-value is not}
Not the probability that $H_0$ is true; not the probability the result is due to chance; not a measure of effect size or importance. A tiny, useless effect can be highly
significant with huge $n$; a large, important effect can be non-significant with small $n$. Always report an effect size and a confidence interval.

\subsection{Errors and power}
\begin{center}\small
\begin{tabular}{l|cc}
& $H_0$ true & $H_0$ false\\\hline
Reject $H_0$ & Type I error (false positive), prob.\ $\alpha$ & correct (power $= 1 - \beta$)\\
Fail to reject & correct & Type II error (false negative), prob.\ $\beta$\\
\end{tabular}
\end{center}
\textbf{Power} rises with sample size, effect size and $\alpha$, and falls with noise. Plan sample size before an experiment (Chapter 28: about $16\sigma^2/\delta^2$ per arm for
80\% power at $\alpha = 0.05$).

\section{The common tests}
\begin{itemize}
\item \textbf{One-sample z/t-test}: is a mean equal to a value? ($t$ when $\sigma$ is estimated, $n - 1$ degrees of freedom.)
\item \textbf{Two-sample t-test}: do two independent groups have equal means? (Welch's version for unequal variances.)
  \begin{example}[: two-sample t]
  Interview scores, group A $(12, 14, 11, 13, 15)$ mean 13; group B $(10, 11, 9, 12, 10)$ mean 10.4. $t = 2.98$ with 8 df, p $= 0.018$: significant difference.
  \end{example}
\item \textbf{Paired t-test} (the syllabus's ``paired means test''): the same units measured twice (before/after a feature, model A vs B on the same folds); test whether the mean
  \emph{difference} is 0.
  \begin{example}[: paired t]
  Five recruiters' time-to-shortlist before and after a tool; differences 2, 1, 3, 0, 4 hours. Mean 2, sd 1.58, SE 0.707, $t = 2.83$ (4 df), p $= 0.047$. Pairing removes
  between-recruiter variation, making the test more powerful.
  \end{example}
\item \textbf{Two-proportion z-test}: do two conversion rates differ? (A/B tests on clicks.)
  \begin{example}[: an A/B test]
  Control: 200 applies out of 4{,}000 (5.0\%); variant: 260 of 4{,}000 (6.5\%). Pooled $\hat p = 460/8000 = 0.0575$; SE $= \sqrt{0.0575\cdot0.9425\cdot(2/4000)} = 0.0052$;
  $z = 0.015/0.0052 = 2.88$; p $= 0.004$. The lift is significant; 95\% CI for the difference $\approx 1.5\pm1.0$ points.
  \end{example}
\item \textbf{Chi-square test of independence}: are two categorical variables related? (Work experience vs placement: $\chi^2 = 16.7$, p $< 0.001$; Chapter 32.)
\item \textbf{ANOVA}: equal means across several groups (Chapter 37).
\item \textbf{Non-parametric}: Mann--Whitney U (two groups, ranks), Wilcoxon signed-rank (paired), Kruskal--Wallis (several groups), Kolmogorov--Smirnov (are two distributions
  the same: drift detection).
\end{itemize}

\section{Multiple testing}
Run 20 independent tests at $\alpha = 0.05$ with no real effects: $P(\text{at least one false positive}) = 1 - 0.95^{20} = 0.64$. Corrections: \textbf{Bonferroni} (use $\alpha/m$:
$0.05/10 = 0.005$ for 10 tests), Holm, and false discovery rate control (Benjamini--Hochberg). Relevant when checking many metrics, segments or features.

\begin{practical}[: where statistics drives ML decisions]
\begin{itemize}
\item \textbf{Model comparison}: a paired test (or bootstrap CI) on per-fold or per-example scores before claiming model B beats model A.
\item \textbf{A/B tests}: two-proportion tests, sample-size planning, guardrails, multiple-testing control; this is how every InfoEdge model launch is judged.
\item \textbf{Uncertainty in metrics}: standard errors of accuracy/AUC on small test sets; bootstrap confidence intervals.
\item \textbf{Drift detection}: KS or chi-square tests comparing feature distributions between training and live data.
\item \textbf{Feature significance} in linear/logistic regression (coefficient t-tests), with heteroscedasticity-robust SEs (Chapter 2).
\end{itemize}
\end{practical}

\stuck{StatQuest: The Central Limit Theorem, Clearly Explained}{\ytid{YAlJCEDH2uY}}
\section{Interview questions}

\begin{qa}{\pyq{reported Round 1: p-value vs alpha} What is a p-value, what is alpha, and how do you compare them?}
\oneline{The p-value is the probability, assuming the null hypothesis is true, of seeing data at least as extreme as observed; alpha is the false-positive rate we decide in advance to tolerate;
if p $\le\alpha$ we reject the null, otherwise we fail to reject it.}
\why Alpha is a property of the decision rule (fixed before the data); the p-value is a property of the observed data. Example: $z = 2$, p $= 0.046$: significant at $\alpha = 0.05$, not at 0.01.
\pushes \emph{``Is p = 0.03 the probability that the null is true?''} No. \emph{``What if p = 0.06?''} Not significant at 5\%; report the effect size and CI, and consider power.
\end{qa}

\begin{qa}{Explain descriptive vs inferential statistics, and population vs sample, with examples.}
\oneline{Descriptive statistics summarise observed data; inferential statistics use a sample to make uncertain statements about a population; a population is the full group of interest (all
seekers), a sample is the observed subset (10{,}000 surveyed seekers), whose statistics estimate population parameters.}
\end{qa}

\begin{qa}{State the central limit theorem and say why it matters for data science.}
\oneline{For independent observations with finite variance, the sample mean is approximately normal with mean $\mu$ and variance $\sigma^2/n$ for large $n$, whatever the original distribution;
it justifies normal-based confidence intervals and tests for means and proportions (A/B tests), and explains why averages are stable.}
\worked Average of 30 die rolls: mean 3.5, SE 0.31.
\end{qa}

\begin{qa}{What is a standard error and how does it differ from a standard deviation?}
\oneline{The standard deviation measures the spread of individual observations; the standard error measures the spread of an estimate (like a sample mean) across repeated samples, $\sigma/\sqrt n$, and
shrinks as $n$ grows.}
\end{qa}

\begin{qa}{Compute a 95\% CI for mean time-to-hire with $n = 36$, $\bar x = 12$, $s = 3$, and interpret it correctly.}
\oneline{$12\pm1.96\times0.5 = (11.02, 12.98)$ days; the procedure captures the true mean in 95\% of repeated samples (not a 95\% probability statement about this fixed interval).}
\end{qa}

\begin{qa}{Type I vs Type II error, and power. How do you increase power?}
\oneline{Type I: rejecting a true null (false positive, rate $\alpha$); Type II: failing to reject a false null (false negative, rate $\beta$); power $= 1 - \beta$; increase it with more data, larger
true effects, less noise (paired designs, variance reduction), or a higher $\alpha$.}
\end{qa}

\begin{qa}{\pyq{syllabus: paired means test} When do you use a paired t-test? Compute it for differences 2, 1, 3, 0, 4.}
\oneline{When the same units are measured twice (before/after, model A vs B on the same data); test whether the mean difference is zero: mean 2, SE 0.707, $t = 2.83$ with 4 df, p $= 0.047$.}
\why Pairing removes unit-to-unit variation, so the test is more powerful than an independent two-sample test on the same numbers.
\end{qa}

\begin{qa}{Analyse an A/B test: control 200/4000 applies, variant 260/4000. Is the lift significant?}
\oneline{Two-proportion z-test: pooled rate 0.0575, SE 0.0052, $z = 2.88$, p $= 0.004$: yes, the 1.5-point lift (30\% relative) is significant at 5\%.}
\pushes \emph{``What else would you check?''} Sample ratio mismatch, guardrail metrics, novelty effects, pre-registered metric, segment heterogeneity, practical significance.
\end{qa}

\begin{qa}{\deep With 50 lakh users in an A/B test, almost everything is significant. How do you decide what matters?}
\oneline{Focus on effect sizes and confidence intervals relative to a pre-defined minimum practically important effect, cost and risk; statistical significance with huge $n$ only says the effect is
not exactly zero.}
\end{qa}

\begin{qa}{What is the multiple testing problem and how do you correct for it?}
\oneline{Testing many hypotheses inflates the chance of at least one false positive ($1 - 0.95^{20} = 0.64$ for 20 tests); correct with Bonferroni ($\alpha/m$), Holm, or Benjamini--Hochberg FDR control.}
\end{qa}

\begin{qa}{\deep How would you test whether model B is really better than model A?}
\oneline{Evaluate both on the same folds or examples and use a paired test (paired t-test on fold scores, McNemar's test on per-example correctness, or a bootstrap CI of the metric difference),
then confirm online with an A/B test.}
\end{qa}

\begin{qa}{Which test would you use: (a) apply rate of two email templates, (b) mean salary in three cities, (c) is city related to job function, (d) has the distribution of experience changed since
training?}
\oneline{(a) two-proportion z-test (or chi-square); (b) one-way ANOVA (Kruskal--Wallis if skewed); (c) chi-square test of independence; (d) Kolmogorov--Smirnov test (or PSI) for drift.}
\end{qa}

\begin{qa}{What is sampling bias? Give an ML example.}
\oneline{A systematic difference between the sample and the population; e.g. training a lead-conversion model only on users that telecallers chose to call, so the model never learns about the users they skipped.}
\end{qa}
